\documentclass[12pt,hyperref=unicode]{beamer}

\usepackage[utf8]{inputenc}
\usepackage[turkish]{babel}
\usepackage[T1]{fontenc}

\usepackage{amsmath}
\usepackage{amssymb}
\usepackage{amsthm}
\usepackage{enumerate}
\usepackage{tikz}
\usepackage{transparent}
\usepackage{xcolor}
\usepackage{listings}
\usepackage{verbatim}
\usepackage{graphicx}

\usepackage{pgfplots}

\usetheme{Warsaw}
\usecolortheme{default}
\definecolor{nigdeyesili_acik}{RGB}{3, 150, 166}
\definecolor{nigdeyesili_koyu}{RGB}{143, 209, 217}
\setbeamercolor{structure}{fg=nigdeyesili_acik}

\author[Dr. Mehmet CANEVİ]{Arş.~Gör.~Dr.~M.~Canevi\inst{1}}

\institute{
    \inst{1}%
    Bilgisayar Mühendisliği\\
    Mühendislik Fakültesi
}
    
\date[2025] {Ders Notları, Ocak 2025}
\setbeamertemplate{footline}[frame number]

\logo{\transparent{0.4}\includegraphics[width = 20mm]{logo}}

\definecolor{mGreen}{rgb}{0,0.6,0}
\definecolor{mGray}{rgb}{0.5,0.5,0.5}
\definecolor{mPurple}{rgb}{0.58,0,0.82}
\definecolor{backgroundColour}{rgb}{0.95,0.95,0.92}

\lstloadlanguages{C}
\lstdefinestyle{CStyle}{
    backgroundcolor=\color{backgroundColour},   
    commentstyle=\color{mGreen},
    keywordstyle=\color{red},
    numberstyle=\tiny\color{mGray},
    stringstyle=\color{mGreen},
    basicstyle=\footnotesize,
    breakatwhitespace=false,         
    breaklines=true,                 
    captionpos=b,                    
    keepspaces=true,                 
    numbers=left,                    
    numbersep=2pt,                  
    showspaces=false,                
    showstringspaces=false,
    showtabs=false,                  
    tabsize=1,
    language=C
}
\lstset{style=CStyle}

\AtBeginDocument{\shorthandoff{=}}
\title[Ders 3] {Karar Verme Yapıları ve Döngüler}
\begin{document}
%%%%%%%%%%%%%%%%%%%%%%%%%%%%%%%%%%%%%%%%%%%%%%%%%%%%%%%%%%%%%%%%%%%%%%%%%%%%%%%%
\frame{\titlepage}
\begin{frame}[fragile]{İçidekiler}
    \tableofcontents
\end{frame}
%%%%%%%%%%%%%%%%%%%%%%%%%%%%%%%%%%%%%%%%%%%%%%%%%%%%%%%%%%%%%%%%%%%%%%%%%%%%%%%%
\section{Karar yapıları}
\begin{frame}[fragile]{if yapısı}
    \lstinline{if} yapısı
    \begin{lstlisting}
        if(sart)
        {
            kod;
        }\end{lstlisting}
\end{frame}
%%%%%%%%%%%%%%%%%%%%%%%%%%%%%%%%%%%%%%%%%%%%%%%%%%%%%%%%%%%%%%%%%%%%%%%%%%%%%%%%
\begin{frame}[fragile]{if-else yapısı}
    \lstinline{if-else} yapısı
    \begin{lstlisting}
        if(sart)
        {
            kod1;
        }
        else
        {
            kod2;
        }\end{lstlisting}
\end{frame}
%%%%%%%%%%%%%%%%%%%%%%%%%%%%%%%%%%%%%%%%%%%%%%%%%%%%%%%%%%%%%%%%%%%%%%%%%%%%%%%%
\begin{frame}[fragile]{if-elseif-else yapısı}
    \lstinline{if-elseif-else} yapısı
    \begin{lstlisting}
        if(sart1)
        {
            kod1;
        }
        else if(sart2)
        {
            kod2;
        }
        else
        {
            kod3;
        }\end{lstlisting}
\end{frame}
%%%%%%%%%%%%%%%%%%%%%%%%%%%%%%%%%%%%%%%%%%%%%%%%%%%%%%%%%%%%%%%%%%%%%%%%%%%%%%%%
\begin{frame}[fragile]{kısa if}
    \lstinline{if-else} yapısını değişken kullanarak kısa ve hızlı bir şekilde yapmak için
    \begin{lstlisting}
        degisken=(sart)? dogru_degeri : yanlis_degeri\end{lstlisting}
    kullanılabilir. Esasen
    \begin{lstlisting}
        if(sart)
        {
            degisken=dogru_degeri;
        }
        else
        {
            degisken=yanlis_degeri;
        }\end{lstlisting}
    koduna karşılık gelir.
\end{frame}
%%%%%%%%%%%%%%%%%%%%%%%%%%%%%%%%%%%%%%%%%%%%%%%%%%%%%%%%%%%%%%%%%%%%%%%%%%%%%%%%
\begin{frame}[fragile]{switch-case yapısı}
    \lstinline{switch-case} yapısı
    \begin{lstlisting}
        switch(degisken)
        {
            case deger1:
              kod1;
              break;
            case deger2:
              kod2;
              break;
            default:
              kod;
        }\end{lstlisting}
\end{frame}
%%%%%%%%%%%%%%%%%%%%%%%%%%%%%%%%%%%%%%%%%%%%%%%%%%%%%%%%%%%%%%%%%%%%%%%%%%%%%%%%
\begin{frame}[fragile]{Uygulama}
    \begin{alertblock}{Soru-1}
        İki sayı ile 4 işlem \lstinline{if} ile yapan programı yazınız.
    \end{alertblock}
    \begin{alertblock}{Soru-2}
        İki sayı ile 4 işlem \lstinline{switch-case} ile yapan programı yazınız.
    \end{alertblock}
\end{frame}
%%%%%%%%%%%%%%%%%%%%%%%%%%%%%%%%%%%%%%%%%%%%%%%%%%%%%%%%%%%%%%%%%%%%%%%%%%%%%%%%
\section{Döngüler}
\begin{frame}[fragile]{while yapısı}
    \lstinline{while} yapısı
    \begin{lstlisting}
        while(sart)
        {
            kod;
        }\end{lstlisting}
\end{frame}
%%%%%%%%%%%%%%%%%%%%%%%%%%%%%%%%%%%%%%%%%%%%%%%%%%%%%%%%%%%%%%%%%%%%%%%%%%%%%%%%
\begin{frame}[fragile]{while yapısı örneği}
    \begin{lstlisting}
        int sayi = 1;

        while(sayi < 5) {
        printf("%d\n", sayi);
        sayi++;
        }\end{lstlisting}
\end{frame}
%%%%%%%%%%%%%%%%%%%%%%%%%%%%%%%%%%%%%%%%%%%%%%%%%%%%%%%%%%%%%%%%%%%%%%%%%%%%%%%%
\begin{frame}[fragile]{for yapısı}
    \lstinline{for} yapısı
    \begin{lstlisting}
        for (ifade1; ifade2; ifade3) {
            kod;
        }\end{lstlisting}
    ve genellikle
    \begin{lstlisting}
        int sayi;
        for (sayi=1; sayi<5; sayi++) {
            printf("%d\n", sayi);
        }\end{lstlisting}
    şeklinde kullanılır;
\end{frame}
%%%%%%%%%%%%%%%%%%%%%%%%%%%%%%%%%%%%%%%%%%%%%%%%%%%%%%%%%%%%%%%%%%%%%%%%%%%%%%%%
\begin{frame}[fragile]{break-continue yapısı}
    \begin{lstlisting}
        int sayi;
        for (sayi=1; sayi<5; sayi++) {
            if(sayi==3)break;
            printf("%d\n", sayi);
        }\end{lstlisting}
        ve
        \begin{lstlisting}
            int sayi;
            for (sayi=1; sayi<5; sayi++) {
                if(sayi==3)continue;
                printf("%d\n", sayi);
            }\end{lstlisting}
        arasındaki fark nedir?
\end{frame}
%%%%%%%%%%%%%%%%%%%%%%%%%%%%%%%%%%%%%%%%%%%%%%%%%%%%%%%%%%%%%%%%%%%%%%%%%%%%%%%%
\end{document}